\section{Format-Specific Reconstruction Methods}\label{sec:format-methods}

\subsection{Common reconstruction and validation requirements}
\label{subsec:format-common}

Each source-to-output profile mapping instantiates
\cref{eq:extended-acceptance} by specifying a constrained output language,
minimum reconstruction level, generated name and media type, resource bounds,
and an authorizing parser reached through a code path separate from the
constructor. This design follows the language-theoretic principle that
acceptance should be defined by an explicit language instead of a catalogue of
malformed cases \cite{momot2016,nail2014}. It also addresses the practical risk
that two file processors assign different boundaries or meanings to the same
bytes \cite{jana2012,ipg2023}.

Complete parsing is stronger than recognition of a prefix. For candidate
\(a=(y,n_o,m_o,k_o,\ell,e)\), the authorizing parser must begin at zero, reach
the unique terminator or declared outer boundary, consume \(|y|\) octets, and
admit every emitted component at its exact parent and position. The output name,
declared media type, byte recognizer, and parser report must identify \(k_o\).
Resource checks use both \(y\) and the recorded reconstruction evidence \(e\).
They cannot be inferred from \(y\) alone. These conditions establish conformity
to the registered output language under the stated assumptions. They do not
establish that arbitrary accepted media are harmless.

The contract for a permitted source profile \(s\), reconstruction level
\(\ell\), and output profile \(k\) is recorded as
\begin{equation}
 \Phi_{p,s,\ell,k}=
 \bigl(\Omega_p,q_{p,s},L_{p,k},\Gamma_{p,k},
       \Pi^{\mathrm{auth}}_{p,v,k},
       \mathcal{N}_{p,k},\mathcal{M}_{p,k},
       \mathcal{B}_{p,s,k}\bigr),
 \label{eq:format-profile-tuple}
\end{equation}
where \(\Gamma_{p,k}\) is the location-aware component and permitted-metadata
policy, \(\mathcal{N}_{p,k}\) and \(\mathcal{M}_{p,k}\) are the generated-name
and output-media-type rules, and \(\mathcal{B}_{p,s,k}\) contains input,
reconstruction, and output bounds. The remaining symbols were defined in
\cref{sec:background,sec:formal-method}. Reconstruction strength follows
\(\NotProcessed\prec\Structural\prec\Semantic\). The instantiated routes and
the limits of their claims are compared in
\cref{tab:format-profile-comparison}.

\begin{table}[H]
\centering
\caption{Comparison of implemented reconstruction routes. ``Retained source''
identifies content that survives a structural route and limits its claim.}
\label{tab:format-profile-comparison}
\scriptsize
\setlength{\tabcolsep}{3.1pt}
\begin{tabularx}{\textwidth}{L{0.15\textwidth}L{0.12\textwidth}L{0.19\textwidth}L{0.21\textwidth}Y}
\toprule
Input profile & Levels & Retained source & Complete-output evidence & Authorized claim \\
\midrule
PNG, JPEG, static WebP & Semantic & Decoded raster only & Grammar parse and second decode & Canonical grammar, exclusion of forbidden source metadata, and source-payload independence \\
GIF & Semantic & Decoded frames and normalized timing & Block parse, trailer closure, and second decode & Animation closure and source-payload independence within the frame model \\
APNG & Structural & \texttt{IDAT}/\texttt{fdAT} payloads & Chunk, CRC, sequence, and frame checks & Container closure, exclusion of forbidden source metadata, and bounded animation structure \\
Animated WebP & Structural & \texttt{VP8}/\texttt{VP8L} frame payloads & RIFF closure, frame checks, and full decode & Container closure, exclusion of forbidden source metadata, and bounded frame consistency \\
MP3 & Structural or semantic & MPEG frames structurally, decoded samples semantically & Exact frame walk or semantic MP3 reparse & Structural tag and suffix removal. Source-payload independence only semantically \\
WAV PCM & Structural or semantic & PCM bytes structurally, decoded samples semantically & RIFF arithmetic or semantic MP3 reparse & Structural container closure. Source-payload independence only semantically \\
FLAC & Structural or semantic & Encoded frames structurally, decoded samples semantically & FLAC frame checks or semantic MP3 reparse & Structural stream closure. Source-payload independence only semantically \\
Ogg Opus/Vorbis & Structural or semantic & Packets structurally, decoded samples semantically & Ogg closure or semantic MP3 reparse & Structural stream closure. Source-payload independence only semantically \\
ISO BMFF & Structural or semantic & Samples in \texttt{mdat} on the structural path & Box grammar, sample tables, and extent coverage & Container and offset closure. Payload independence only after semantic reconstruction \\
WebM & Structural or semantic & Blocks structurally, decoded streams semantically & EBML closure or semantic MP4 reparse & Structural container claims. Source-payload independence only semantically \\
\bottomrule
\end{tabularx}
\end{table}

Policy fixes \(\Phi_{p,s,\ell,k}\) before dispatch. The handler checks
\(\mathcal{B}_{p,s,k}\), reconstructs the object, and reports the level
performed. The authorizing parser then evaluates the complete candidate against
\(L_{p,k}\) and \(\Gamma_{p,k}\) before rebuilt-output scanning and publication.
Overflow, an incomplete boundary, a forbidden component, an exhausted budget,
or \(\ell<q_{p,s}\) prevents delivery.

The strict policy maps supported audio sources to an MP3 output profile and
supported video sources to an MP4 output profile after semantic reconstruction.
The source profile and delivered profile therefore need not be identical. The
reference encoder, concrete codec settings, and authorizing-parser identifiers
are implementation parameters reported in \cref{sec:implementation}. A WebM
label in the evaluation denotes the source profile. Its semantic route emits an
MP4 candidate.

Component closure and metadata policy are related but not identical. Closure
defines the complete set of legal structural units and their permitted
locations. Metadata policy determines which legal but nonessential units may
survive. The evaluated profiles use removal or replacement rather than
source-field copying. A byte-substring scan is not used to establish removal,
because an arbitrary byte sequence can occur inside compressed data. Absence
is established by reparsing the output grammar and observing that no forbidden
component is present.

Size and offset integrity are also treated as parser obligations. For a
length-delimited component $c$ at offset $o_c$ with header length $h_c$ and
declared payload length $\ell_c$, the parser requires checked arithmetic and
the containment relation
\begin{equation}
o_c+h_c+\ell_c\leq b_{\mathrm{parent}},
\qquad
[o_c,o_c+h_c+\ell_c)\subseteq B_{\mathrm{parent}}.
\label{eq:component-bounds}
\end{equation}
For sequential grammars, the next component must begin exactly at the previous
component's end. For offset-bearing grammars, a referenced extent must lie in
an allowed media-data region and must agree with the relevant count and sample
tables. Checks precede slicing, allocation, or offset conversion. These rules
turn complete consumption into a constructive invariant rather than an
end-of-function comparison alone.

\subsection{Static raster images}\label{subsec:static-raster-method}

JPEG, PNG, static WebP, BMP, and TIFF are accepted as static raster input
classes when the image library can probe and decode them under policy limits.
The output profiles specialize the corresponding JPEG, PNG, and WebP
specifications rather than treating successful decoding as the grammar
definition~\cite{ituT81,png3,webpRiff}.
Only PNG, JPEG, and static WebP are output classes. The handler first probes
dimensions without decoding the full pixel buffer. Zero dimensions, probe
failure, or a probed pixel count above the policy limit block before full
decode. After decode, dimensions and pixel count are checked again. If the
decoded pixel count exceeds twice the probed count, the operation is rejected.
Any other nonzero discrepancy produces a warning and therefore remains
non-deliverable under enforcement mode, while both counts remain subject to
their respective policy bounds. The intermediate
state is a decoded raster, so the resulting method is semantic reconstruction.
BMP and TIFF source syntax, embedded profiles, and application records do not
survive because those containers are not copied into the output.

The selected encoder creates a new raster object from decoded pixel state.
The implementation's image representation does not carry source EXIF, XMP,
ICC, comment, or thumbnail records through this path. The authorizing parser
then performs both structural parsing and a second decode. The second probe
and decode must agree on dimensions and remain below the validator's pixel
bound. This post-reconstruction decode matters because structural
correctness of marker or chunk boundaries alone does not establish that the
generated compressed payload is consumable by the selected decoder.

For PNG, the accepted grammar is
\begin{equation}
\mathrm{Sig}\parallel\mathrm{IHDR}\parallel[\mathrm{PLTE}]
\parallel[\mathrm{tRNS}]\parallel\mathrm{IDAT}^{+}
\parallel\mathrm{IEND}.
\label{eq:png-grammar}
\end{equation}
The parser verifies the PNG signature, every chunk length, every CRC, the
single 13-byte \texttt{IHDR}, color-type constraints, palette and transparency
compatibility, contiguous nonempty image-data chunks, and an empty terminal
\texttt{IEND}. No ancillary chunk other than the profile-permitted
\texttt{tRNS} is admitted, and unknown chunks are rejected. Animation chunks
are forbidden in the static profile. The terminal offset must equal file
length. These constraints implement a narrower language than
the full PNG specification \cite{png3}.

For JPEG, the output begins with SOI and one pinned generated JFIF APP0 segment.
Its identifier, version, density unit, nonzero density fields, and zero
thumbnail dimensions are checked. The validator uses an explicit marker
allowlist and admits baseline
SOF0 and progressive SOF2 output. It checks frame dimensions and component
table lengths, requires valid scan headers, walks entropy-coded data with JPEG
stuffing and restart-marker rules, and ends at one EOI. Other APP segments and
COM are forbidden. Marker lengths must remain within the file and no byte may
follow EOI. The profile therefore excludes source
application segments rather than attempting to distinguish benign from
dangerous values inside them.

For static WebP, the validator treats the file as a RIFF form. It requires the
outer size to satisfy $\mathrm{RIFFSize}=\lvert y\rvert-8$, the form type
\texttt{WEBP}, exact chunk boundaries, and the specified even-byte padding.
Exactly one \texttt{VP8 } or \texttt{VP8L} image payload is present. When
\texttt{VP8X} is emitted, its feature flags, dimensions, animation state, and
alpha declaration must agree with the present chunks. Animation and metadata
chunks, including \texttt{ANIM}, \texttt{ANMF}, \texttt{EXIF}, \texttt{XMP },
and \texttt{ICCP}, are excluded. A second WebP decode supplies the semantic
postcondition.

\subsection{GIF, APNG, and animated WebP}\label{subsec:animated-raster-method}

GIF uses semantic frame reconstruction under a constrained specialization of
the GIF89a block grammar~\cite{gif89a}. The handler decodes frames to RGBA,
checks the configured frame count and per-frame pixel limit, and accumulates a
separate total-pixel budget. Frame delays below two centiseconds are
normalized to two centiseconds. The sum of frame delays for one animation cycle
is bounded by 60 seconds. This bound does not limit total viewer time when the
generated loop declaration requests indefinite repetition. Re-encoding
constructs a fresh GIF with generated color tables and that explicit loop
declaration. Source
comment, plaintext, and application-extension data are not copied. Frame
offsets, dimensions, disposal operation, transparency, and normalized delay
are represented in generated frame records.

The authorizing parser accepts GIF87a and GIF89a signatures, requires one
graphic-control block per frame, and permits only the generated
\texttt{NETSCAPE2.0} loop extension. It checks the logical screen descriptor,
color tables, image rectangles, Lempel--Ziv--Welch (LZW) sub-block chains, and
one terminal trailer. Unknown extensions and trailing bytes fail validation. A
second decoder must read at least one frame and reach end of input. The current
parser does not couple extension use to the GIF89a signature, so that
version-extension check remains a release condition. The semantic boundary is
decoded frame content under the stated delay normalization. Pixel fidelity is
evaluated separately.

APNG differs because the implemented handler is a structural chunk remux. It
walks the PNG chunk sequence, validates chunk boundaries and CRCs, retains the
compressed \texttt{IDAT} and \texttt{fdAT} payloads, and removes ancillary
chunks not needed by the constrained animation. The output grammar is
\begin{equation}
\mathrm{Sig}\parallel\mathrm{IHDR}\parallel[\mathrm{PLTE}]
\parallel[\mathrm{tRNS}]\parallel\mathrm{acTL}
\parallel\mathrm{Frames}\parallel\mathrm{IEND}.
\label{eq:apng-grammar}
\end{equation}
There is one \texttt{acTL} before image data. The default \texttt{IDAT} image
may be the first animation frame when preceded by its \texttt{fcTL}, or it may
remain outside the animation when the first \texttt{fcTL} follows
\texttt{IDAT}. Each \texttt{fcTL} describes a nonzero rectangle within the
canvas and uses defined disposal and blend operations. The shared
\texttt{fcTL}/\texttt{fdAT} sequence starts at zero and is gap-free and
duplicate-free. Every controlled frame has data, and the parsed frame count
equals the \texttt{acTL} declaration. Only \texttt{IHDR}, conditionally required
\texttt{PLTE} and \texttt{tRNS}, \texttt{acTL}, \texttt{fcTL}, \texttt{IDAT},
\texttt{fdAT}, and \texttt{IEND} are allowed.

Because \texttt{IDAT} and \texttt{fdAT} payloads survive, APNG instantiates the
structural claim in \cref{tab:format-profile-comparison}. A policy that requires
semantic reconstruction blocks this route.

Animated WebP has a separate structural profile rather than being admitted by
the static WebP parser. The handler requires an exact RIFF boundary, one
10-byte \texttt{VP8X} first with the animation feature set, one \texttt{ANIM}
before all frames, and at least one \texttt{ANMF}. Canvas dimensions, frame
rectangles, frame count, per-frame pixels, and aggregate pixels are bounded.
Each frame contains either an optional \texttt{ALPH} followed by one nonempty
\texttt{VP8 } payload, or one nonempty \texttt{VP8L} payload without a separate
\texttt{ALPH} chunk. The frame alpha state and the \texttt{VP8X} alpha flag must
agree. Chunk sizes and even-byte padding are consumed exactly. The handler removes \texttt{ICCP},
\texttt{EXIF}, and \texttt{XMP } chunks, clears their \texttt{VP8X} feature
bits, rejects every other top-level component, and regenerates the RIFF size.

The \texttt{webp-animated-canonical-v1} validator, invoked separately from the
constructor, repeats the complete structural parse and additionally checks
encoded frame dimensions, alpha-state consistency, and agreement between
parsed and decoded frame counts. It decodes every emitted frame under fixed
frame and aggregate-pixel bounds and verifies the declared canvas. This
corroborates decoder consumption but remains structural because the source
\texttt{VP8 } or \texttt{VP8L} payloads survive. A policy requiring semantic
reconstruction therefore rejects this route.

\subsection{MP3, WAV, and FLAC}\label{subsec:native-audio-method}

The MP3 structural path recognizes and removes an initial ID3v2 region and a
terminal ID3v1 tag, then validates the remaining bytes as a nonempty sequence
of length-accounted MPEG audio frame extents with admissible
headers~\cite{iso11172}. The frame parser checks sync,
MPEG version,
layer, bitrate index, sampling-rate index, and padding. It computes each frame
length from the header and requires the next header at exactly that boundary.
The canonical output profile admits no ID3, APE, Lyrics3, arbitrary padding,
inter-frame material, or suffix. Stream parameters are required to remain
stable under the implemented profile. Retained MPEG frames limit this path to
the structural MP3 claim in \cref{tab:format-profile-comparison}.

WAV is reconstructed as a new RIFF/WAVE container around retained sample
bytes under the RIFF interchange model~\cite{microsoftRiff}. The input walk
checks RIFF size, chunk boundaries, and padding while
locating one supported \texttt{fmt } chunk and one \texttt{data} chunk.
Source ancillary chunks are omitted. The output uses a 16-byte format body for
integer PCM, followed by the data chunk and generated zero padding when needed.
The validator requires exactly that order and no other chunks. It checks
nonzero channel count and sample rate, supported sample width, block
alignment, byte rate, and data divisibility. These conditions require
\begin{equation}
\begin{aligned}
\mathrm{blockAlign}&=C\frac{B}{8},\\
\mathrm{byteRate}&=F_s\,\mathrm{blockAlign},\\
\lvert\mathrm{data}\rvert\bmod\mathrm{blockAlign}&=0,
\end{aligned}
\label{eq:wav-integrity}
\end{equation}
where $C$ is channel count, $B$ is bits per sample, and $F_s$ is sample rate.
The RIFF size includes any generated pad byte and equals $\lvert y\rvert-8$.
Although PCM sample bytes are retained, there is no nested compressed codec
bitstream in this constrained profile. The reconstruction is still reported
as structural because sample semantics are not decoded and regenerated.

FLAC structural reconstruction follows the native stream organization
standardized in RFC~9639~\cite{rfc9639}. It retains the 34-byte STREAMINFO payload and
encoded audio frames while removing every other metadata block. It emits the
\texttt{fLaC} marker, exactly one STREAMINFO block marked as the final metadata
block, and one or more retained frames. VORBIS\_COMMENT, PICTURE, APPLICATION,
SEEKTABLE, CUESHEET, PADDING, and unknown metadata blocks are absent from the
output grammar.

The separate FLAC authorizing parser is deeper than a metadata-chain walk. It
checks STREAMINFO block-size, frame-size, sample-rate, channel, bit-depth, and
total-sample fields with checked arithmetic. For each frame it parses the sync
and blocking strategy, canonical coded frame or sample number, block-size and
sample-rate codes, channel assignment, explicit bit depth, header CRC-8, and
frame CRC-16. It traverses constant, verbatim, fixed-predictor, and LPC
subframes, including wasted-bit declarations, Rice and escaped residual
partitions, and zero alignment. Frame numbering, aggregate sample count, and
parameters are cross-checked with STREAMINFO. Streamable-subset limits are
applied to block size, low-rate LPC order, and Rice partition order. Every
frame end becomes the next frame start, and the final frame ends at EOF.

These checks establish a constrained, frame-CRC-consistent FLAC grammar.
Predictor output and STREAMINFO's MD5 over decoded samples are not recomputed,
so retained frames remain outside payload-neutralization claims.

\subsection{Ogg Opus and Vorbis}\label{subsec:ogg-method}

The implemented Ogg method accepts one logical Opus or Vorbis stream under the
page model defined by RFC~3533~\cite{rfc3533}. The
handler validates page capture patterns, version, lacing boundaries, CRC,
serial number, sequence continuity, and BOS/EOS state. It requires a
recognized identification packet, a replaceable comment packet, subsequent
media, and a terminal EOS page. Reconstruction replaces the comment packet
with a minimal generated OpusTags or Vorbis comment. The current replacement
is limited to a comment packet wholly contained in one page and
not sharing that page with another packet. Inputs that would require a more
general repagination are blocked rather than approximated.

The authorizing parser reconstructs packets across page lacing
tables. Continuation flags must agree with whether the prior page ended in an
unfinished packet. Exactly one serial number is allowed, page sequence
numbers are contiguous, every page CRC is valid, and EOS ends the stream at
EOF. Codec header and comment order are checked against the selected Opus or
Vorbis grammar. The Vorbis packet obligations follow the codec specification,
while Opus follows the Ogg framing and header profile in RFC~7845
\cite{vorbisI,rfc7845}. Encoded media packets remain structural evidence only.
The authorized claim is stated in \cref{tab:format-profile-comparison}.

\subsection{ISO BMFF}\label{subsec:isobmff-method}

MP4, MOV, and M4A inputs enter a common ISO base media file format path. The
profile is a constrained use of the ISO BMFF box model and its registered
code points~\cite{iso1449612,mp4ra}. The
native handler performs structural rewriting. It removes excluded metadata
and free-space structures, recursively rebuilds permitted boxes, normalizes
selected timestamps, language fields, handler names, and visual compressor
strings, and retains samples in \texttt{mdat}. The implementation bounds
nesting depth and atoms per level. Because removal changes layout, chunk
offsets are adjusted only after the transformed top-level arrangement is
known. The current native fix-up uses a checked displacement for collected
\texttt{stco} and \texttt{co64} tables. For audio routing, the media-type
aliases \texttt{audio/mp4}, \texttt{audio/x-m4a}, and \texttt{audio/m4a}
converge on the same M4A input profile. This alias convergence is necessary
for recordings produced by common mobile media APIs, but it does not define a
new output grammar. The authorizing parser
accepts a narrower nonfragmented, single-\texttt{mdat} grammar so that this
offset model can be checked end to end.

The canonical hierarchy contains one profile-allowed \texttt{ftyp} brand set,
one structurally
complete \texttt{moov}, and one \texttt{mdat}. Box sizes may use standard or
extended forms only when their declared end lies within the enclosing parent.
Unknown and forbidden children fail according to parent location. The
validator excludes \texttt{udta}, metadata-bearing \texttt{meta},
\texttt{uuid}, free space, attachments, arbitrary vendor boxes, and bytes
outside top-level boxes. The structural rewriter also removes records excluded
by the constrained profile, including \texttt{sgpd}, \texttt{sbgp},
\texttt{sdtp}, \texttt{stps}, \texttt{cslg}, \texttt{iods}, \texttt{elng}, and
\texttt{tref}. Some of these boxes can affect sample grouping, dependency,
presentation, seeking, or inter-track relationships. Their removal therefore
supports a narrower container-closure claim, not semantic equivalence or
general playback compatibility. Structural delivery is confined to registered
source profiles for which the policy permits this loss. Otherwise semantic
transcoding is required. Compatibility remains an empirical fidelity outcome.
Permitted
\texttt{fiel} and \texttt{chrm} sample-entry children are accepted only with
their constrained field values and lengths. The validator checks mandatory children of movie, track, media,
media-information, and sample-table boxes. Data references must be
self-contained. Handler type, sample-entry codec, dimensions, channel fields,
timescales, durations, and edit lists are constrained.

Sample-table integrity is established jointly rather than box by box.
Time-to-sample counts, composition counts when present, sample-size counts,
sync-sample indices, sample-to-chunk runs, and chunk-offset counts must describe
one coherent sample sequence. For every chunk $j$, the validator derives the
number and size of its samples and requires
\begin{equation}
[o_j,o_j+L_j)\subseteq B_{\mathtt{mdat}},
\label{eq:isobmff-offset}
\end{equation}
where $o_j$ comes from \texttt{stco} or \texttt{co64} and $L_j$ is derived from
\texttt{stsc} and \texttt{stsz}. Fragment boxes are outside the current output
grammar. Across all tracks, the derived extents are sorted and must
cover the single \texttt{mdat} payload exactly, without a gap or overlap. A
semantic transcoding path may produce an MP4 candidate, but its output is
normalized by the structural rewriter and must pass the same authorizing parser
before publication.

The structural path establishes container closure and offset consistency for
the constrained nonfragmented profile. Payload independence requires the
semantic path. Strict and paranoid policies block when it is unavailable.

\subsection{WebM}\label{subsec:webm-method}

The implemented EBML output profile is WebM only. Matroska is not promoted
to support by shared EBML recognition. The profile is defined against the EBML
and Matroska element model while admitting only the WebM subset named by the
output grammar~\cite{rfc8794,rfc9559,webmContainer,matroskaElements}. The structural handler requires one
bounded EBML header with \texttt{DocType=webm}, one Segment, one Info element,
one Tracks element, and at least one Cluster. It removes SeekHead, Cues, Tags,
Attachments, Chapters, Void, checksums that become stale, and cluster-relative
Position and PrevSize records. Info is rebuilt in a fixed order. Source
application names, date, title, and unsupported fields are removed, while
MuxingApp and WritingApp are set to a fixed implementation identifier. Segment
and child sizes are regenerated with checked EBML variable-length integers.

The \texttt{webm-constrained-v2} validator reparses the full nested grammar on a
code path separate from construction. Element identifiers must use their
shortest valid VINT encoding.
Declared sizes must consume their parents exactly. Nonminimal size widths are
accepted because EBML permits them, while unknown size is allowed only for the
root Segment whose end is fixed by EOF. The accepted Segment orders Info,
Tracks, and Clusters and excludes unknown children.

The track profile permits one VP8, VP9, or AV1 video track and at most one Opus
audio track. Track numbers and unique identifiers are distinct. Dimensions,
pixel count, duration, timestamp scale, sample rate, and channel count are
bounded. The validator checks VP9 feature records, AV1 configuration-record
shape and OBU boundaries, and the complete Opus identification header,
including mapping, pre-skip-derived CodecDelay, and SeekPreRoll. Agreement
between AV1 configuration fields and an optional Sequence Header OBU remains a
codec-semantic obligation and is not claimed.

Each Cluster begins with one Timestamp and contains only SimpleBlock media.
Block track numbers must identify declared tracks. Reserved flags and all
supported lacing modes are parsed to positive frame extents. Video lacing is
excluded. Cluster, block, frame, aggregate frame-count, timestamp, and
one-hour duration bounds are enforced. Absolute block timestamps are
nondecreasing across Cluster boundaries, declared Duration cannot end before
the final block timestamp, and each declared track must carry a frame. The
last parsed Segment byte must be the last file byte.

The native WebM validator closes the container grammar without decoding VP8,
VP9, AV1, or Opus frames. A structural WebM candidate therefore cannot satisfy a
policy that requires semantic reconstruction. The strict route decodes the
source and emits validated MP4 rather than claiming semantic WebM output.

\subsection{Post-reconstruction validation and deferred evidence}
\label{subsec:post-reconstruction-validation}

Every method ends at the predicate in \cref{eq:extended-acceptance}. The
authorizing report, produced separately from construction, must identify the
expected profile and place its terminal offset at $|y|$ before rebuilt-output
signature policy is evaluated. Decoder exit, MIME inference, handler return,
and signature scanning cannot substitute for complete grammar validation.

The implementation and unit tests identify where these checks are enforced.
Empirical claims are limited by the executed protocol in \cref{sec:evaluation}.
Incomplete runs are reported only as finite observations or negative results.
